% aRES N-block charge architecture -- short findings note.
% Build:  latexmk -pdf charge_architecture.tex

\documentclass[11pt,a4paper]{article}

\usepackage[T1]{fontenc}
\usepackage[utf8]{inputenc}
\usepackage[margin=2.4cm]{geometry}
\usepackage{graphicx}
\usepackage{booktabs}
\usepackage{enumitem}
\usepackage[colorlinks=true,linkcolor=blue,urlcolor=blue]{hyperref}

\setlength{\parindent}{0pt}
\setlength{\parskip}{6pt}

\title{\vspace{-1.3cm}\textbf{A bipolar charge architecture in the aRES N-terminal block}}
\author{\normalsize Rotation project \texttt{aRES\_ppi}}
\date{\normalsize\today}

\begin{document}
\maketitle
\vspace{-0.7cm}

Sequence composition and AlphaFold3 model analysis only --- no experimental data.
\emph{E.~coli} UC5 aRES, 459 aa.

\section*{1. The N-terminal block is two modules of opposite charge}

Per-region charge composition, counting Lys$+$Arg against Asp$+$Glu:

\begin{center}
\begin{tabular}{lrrrrrr}
\toprule
Region & Length & K & R & D & E & Net (per 100 aa) \\
\midrule
Zn ribbon 1--48      &  48 & 10 &  3 &  2 &  0 & $+11$ \ \ ($+22.9$) \\
NTD-1 49--120        &  72 &  0 &  3 & 11 &  8 & $-16$ \ \ ($-22.2$) \\
N-block 1--120       & 120 & 10 &  6 & 13 &  8 & $-5$ \ \ ($-4.2$) \\
RES side 126--447    & 322 & 23 & 19 & 24 & 21 & $-3$ \ \ ($-0.9$) \\
\bottomrule
\end{tabular}
\end{center}

Two directly adjacent modules carry almost exactly equal and opposite charge density. The zinc
ribbon has ten lysines and \emph{no glutamate}; NTD-1 has nineteen acidic residues and \emph{no
lysine at all}. The RES domain next door is at $-0.9$ per 100 residues, i.e.\ an ordinary protein
surface --- so the extremity is confined to the N-block and is anticorrelated between its halves.

Residues 31--40 are \texttt{CSYC\,KRKRK\,N}: a five-residue basic run immediately C-terminal to the
second CxxC motif. That is the canonical arrangement of a nucleic-acid-binding zinc ribbon, and it
is a stronger signal than the net charge.

\section*{2. Both charged faces are solvent-exposed}

In the electrostatic surface of the monomer model, the basic ribbon face (carrying the zinc) and the
acidic NTD-1 face sit on adjacent lobes and each faces solvent. They do not neutralise one another
in an intramolecular clamp. The N-block therefore presents a \textbf{bipolar surface}: one strongly
basic face and one strongly acidic face, both available for binding.

\section*{3. Working model --- hypothesis, not a result}

The N-block may recognise \textbf{a polymerase engaged with nucleic acid}, with the two modules
reading the two components:

\begin{itemize}[leftmargin=1.4em,itemsep=2pt]
  \item the \textbf{basic zinc ribbon} binds the DNA;
  \item the \textbf{acidic NTD-1} binds the polymerase's basic nucleic-acid-binding cleft, acting as
        a DNA mimic (the logic of Ocr, T7's own gene 0.3 product, which mimics B-form DNA to inhibit
        host EcoKI; protein mimics of DNA are acidic).
\end{itemize}

What recommends it:

\begin{itemize}[leftmargin=1.4em,itemsep=2pt]
  \item \textbf{Promiscuity.} Every polymerase has a basic DNA-binding cleft regardless of sequence
        family. aRES must fire on T7 gp5, on divergent T-even and T5 polymerases, and on something in
        $\lambda$. Recognising a conserved \emph{physical} feature does not require a conserved epitope.
  \item \textbf{Specificity.} Requiring both faces to engage is an AND-gate --- which is what the
        system needs, since the host cytoplasm is full of polymerases and a false positive is suicide.
  \item \textbf{Why the sensor is bimodular at all}, which neither a folding-scaffold nor an
        effector role for the ribbon explains.
  \item \textbf{$\lambda$.} Host Pol III hijacked at ori$\lambda$ by P/O is still a polymerase on DNA
        in an aberrant context, consistent with the host-factor-hijacking model.
  \item \textbf{It survives the DNA-free reconstitution.} Purified gp5$\cdot$TrxA activates aRES with
        no DNA and no dNTPs. If NTD-1 is the DNA mimic it engages the cleft without DNA present; the
        ribbon's DNA arm would contribute avidity and context-specificity \emph{in vivo}. A sufficient
        minimal reaction and the full mechanism need not be the same thing.
\end{itemize}

\section*{4. Tests, cheapest first}

\begin{enumerate}[leftmargin=1.6em,itemsep=2pt]
  \item \textbf{Artifact control on the electrostatics.} Recompute with the Zn deleted, and colour the
        surface by residue type as a model-free check
        (\texttt{analysis/scripts/color\_by\_restype.pml}). The blue is currently centred on the metal,
        and if the four Cys were not treated as thiolates the site is over-positive by up to $+6$.
  \item \textbf{Back face and contiguity.} If the two patches form one contiguous bipolar surface, the
        two-sided model is viable; if they point away from each other they cannot engage one target at
        once and the model needs rethinking.
  \item \textbf{Is D104 on the acidic face?} The existing cofold places aRES D104 against gp5 S122. If
        D104 sits in the exposed acidic patch, the model connects to the one interface result in hand,
        and the D104A/D104N experiment tests the model rather than a single weak contact.
  \item \textbf{AF3 cofold of aRES with a primer--template duplex.} Direct test: does DNA dock onto the
        KRKRK face? One server job.
  \item \textbf{DNA binding, 1--120 vs 49--136.} EMSA or fluorescence anisotropy against ssDNA, dsDNA
        and a primer--template fork. 49--136 contains no ribbon and no zinc, so it is the built-in
        negative control and is already in the construct series.
  \item \textbf{DNA competition of the aRES--gp5 interaction.} If NTD-1 occupies the DNA cleft, adding
        primer--template should compete aRES off gp5. Sharpest available test of the mimicry half.
  \item \textbf{Salt dependence of $K_d$.} An electrostatically driven interface is strongly
        ionic-strength sensitive; a conventional epitope interface much less so.
  \item \textbf{Conservation.} Compute the same per-module charge split across PF18870 homologs. If
        ribbon-containing homologs consistently show the mirror, it is functional architecture rather
        than an accident of one sequence.
\end{enumerate}

\section*{5. Caveats}

\begin{itemize}[leftmargin=1.4em,itemsep=2pt]
  \item This is compositional reasoning over a single predicted model. It generates a hypothesis; it
        does not confirm one.
  \item The electrostatic rendering may overstate the basic patch at the zinc (test 1).
  \item The surface was inspected in one orientation (test 2).
  \item gp5 D106 is itself acidic, which is awkward for a naively electrostatic interface. Local
        context will matter more than net charge.
  \item The DNA-free \emph{in vitro} activation bounds any model in which nucleic acid is the trigger
        for T7. It does not exclude a contribution to affinity or to other triggers.
\end{itemize}

\section*{6. Figures to make}

\begin{enumerate}[leftmargin=1.6em,itemsep=2pt]
  \item \textbf{\texttt{figs/restype\_front.png}, \texttt{figs/restype\_back.png}} --- surface coloured
        by residue type, front and back of the N-block. Script:
        \texttt{analysis/scripts/color\_by\_restype.pml}.
  \item \textbf{\texttt{figs/apbs\_nozn.png}} --- the existing electrostatic surface recomputed with the
        Zn removed, same orientation and same $\pm 5\,k_{\mathrm{B}}T/e$ scale, beside the original.
  \item \textbf{\texttt{figs/charge\_window.png}} --- net charge in a 15-residue sliding window along
        residues 1--459, with the module boundaries marked. Shows the anticorrelation at 48 as one line.
  \item \textbf{\texttt{figs/nblock\_cartoon.png}} --- cartoon of 1--120 with the ribbon, NTD-1 and the
        zinc coloured, oriented to match figure 1.
\end{enumerate}

\end{document}
